\documentclass{article}

\usepackage[utf8]{inputenc}     % pdfLaTeX
\usepackage[T1]{fontenc}
\usepackage[magyar]{babel}
\usepackage{amsmath}

% Set page size and margins
% Replace `letterpaper' with `a4paper' for UK/EU standard size
\usepackage[a4paper,top=2cm,bottom=2cm,left=3cm,right=3cm,marginparwidth=1.75cm]{geometry}
%for arrows
\usepackage{textcomp}

\title{Márkanevek Helyesírása}
\author{Lengyel Zoltán Péter \thanks{Csukás Ibolya tanítása alapján}}
\date{Legutóbb frissítve: 2026. Március 16.}

\begin{document}

\maketitle

\tableofcontents

\newpage

\begin{itemize}
    \item Nagybetűvel, de csak a márkanevet
    \begin{itemize}
        \item pl. Milka csoki, Tomi Kristály mosópor, Túró Rudi, Egri Bikavér
    \end{itemize}
    \item Kiegészülhet számokkal, betűkkel
    \begin{itemize}
        \item pl. Lewis 501, Hélia-D
    \end{itemize}
\end{itemize}

\section{Toldalékolás}
\begin{itemize}
    \item Toldalékolás: közvetlenül
    \begin{itemize}
        \item pl. Adidast, Túró Rudit
        \item de! tőnyújtás! \begin{itemize}
            \item pl. Alda Rómeót, Coca-Colát
        \end{itemize}
    \end{itemize}
    \item Hosszú mássalhangzóra végződő márkanév esetén a -val/-vel hasonult ragot kötőjellel kapcsoljuk \begin{itemize}
        \item pl. Knorr-ral, Tomi Matt-tal
    \end{itemize}
    \item Néma hangra végződő márkanevekhez kötőjellel kapcsoljuk a kiejtés szerinti ragot
    \begin{itemize}
        \item Peugeot-ról, Renault-val
    \end{itemize}
\end{itemize}

\section{-i, -s képzős alakok}
\begin{itemize}
    \item Egytagú máranév esetében kisbetűvel írjuk és a képzőt közvetlenül csatoljuk (kivéve néma hangra, vagy hosszú mássalhangzóra végződő márkaneveknél, ott kötőjellel) \begin{itemize}
        \item pl. adidasos, nike-s, colgate-es
    \end{itemize}
    \item Kéttagú (vagy többtagú) márkanév esetén megtartjuk a nagybetűket és a toldalékot kötőjellel kapcsoljuk \begin{itemize}
        \item pl. Coca-Colá-s, Alfa Romeó-s
    \end{itemize}
\end{itemize}

\end{document}